\documentclass[10pt,a4paper]{amsart}
\usepackage[margin=19mm]{geometry}
\usepackage{amsmath,amssymb,mathtools}
\usepackage[scr=rsfs,cal=euler]{mathalfa}
\usepackage{xfrac}
\usepackage[unicode,hidelinks,bookmarks=false]{hyperref}
\newcommand{\ie}{i.e.\ }
\newcommand{\cf}{cf.\ }
\newcommand{\resp}{respectively\ }
\newcommand{\Subsection}{\subsection}
\providecommand{\nobreakdash}{\nobreak\hbox{-}}
\setlength{\emergencystretch}{2em}
\multlinegap0pt
\usepackage{bm}
\usepackage{bbm}
\usepackage{dsfont}
\usepackage{xspace}
\usepackage{cancel}
\usepackage{mathtools}
\usepackage{amsthm}
\makeatletter
\@ifundefined{theorem}{\newtheorem{theorem}{Theorem}}{}
\makeatother
\title[Non-local Boundary Value Problems, stochastic resetting and ]{Non-local Boundary Value Problems, stochastic resetting and Brownian motions on graphs}
\author{Stefano Bonaccorsi, Fausto Colantoni, Mirko D'Ovidio, Gianni Pagnini}
\date{Source: arXiv:2209.14135v4 (CC BY 4.0)}
\begin{document}
\maketitle

\noindent\emph{Source and licence.} Non-local Boundary Value Problems, stochastic resetting and Brownian motions on graphs. Version arXiv:2209.14135v4 (CC BY 4.0), distributed under the Creative Commons Attribution 4.0 International licence, as recorded in the arXiv licence metadata for this exact version and shown on the abstract page https://arxiv.org/abs/2209.14135v4. Full rights notice: licenses/arxiv-2209.14135-CC-BY-4.0.txt.\par\smallskip

\noindent\textit{Editorial scope.} Counted unit: the Theorem whose source label is \texttt{thm:Xpallinodelay} in the article (source0.tex lines 1268--1282, source label \texttt{thm:Xpallinodelay}), delivered together with the article's own complete original proof of it (source0.tex lines 1283--1326, one \texttt{proof} environment of 44 lines). No mathematics is written, repaired or re-derived: statement and proof are the article's own text line for line. Required context retained uncounted: fracdyntm (lines 470--485); upsilonSol (lines 1436--1439); timeBondCond (lines 1516--1519); I1piuI2 (lines 1608--1613); resolvent2 (lines 1666--1669); marchaud-graphs (lines 1724--1728). Every definition, display and label of those lines is retained unchanged. Omitted from this package and not redistributed: 1--469, 486--1267, 1327--1435, 1440--1515, 1520--1607, 1614--1665, 1670--1723, 1729--2264. Nothing inside a retained excerpt is omitted or altered: every retained body is delivered exactly as the article's lines are written, so no omission receipt of any kind accompanies this package. Citations: the retained excerpts carry no citation command at all, so no bibliography entry of the article is transcribed and no reference is redistributed. Reference adaptation: none was needed and none was made; every reference inside the delivered unit resolves inside it, so no unresolved reference or citation is produced. Printed slips kept literal: no printed word, symbol, hypothesis or number of the counted statement or of its proof was altered. Layout and class: the article's own class is \texttt{article} with packages including fontenc, inputenc, lmodern, amsmath, amssymb, amsfonts, amsthm, mathrsfs, graphicx, xcolor, hyperref, authblk, bm, enumerate; the wrapper is the lane's pinned \texttt{amsart} editorial wrapper at 10pt on A4, which restates the theorem-like environments the retained text needs and adds the article's own macro definitions that the retained text needs (none), with extra wrapper packages (bm, bbm, dsfont, xspace, cancel, mathtools). The delivered numbering is the unit's own. Assets: no retained span contains a figure environment, a graphics-inclusion command, a PDF-page inclusion, a table or a TikZ picture; no image file is copied into this package, the package declares no asset dependency and no image file is redistributed with it. Licence: version arXiv:2209.14135v4 of this article is distributed under the Creative Commons Attribution 4.0 International licence, as recorded in the arXiv licence metadata for this exact version and shown on the abstract page https://arxiv.org/abs/2209.14135v4. A full rights notice is delivered at licenses/arxiv-2209.14135-CC-BY-4.0.txt. Characters: every retained excerpt is pure ASCII, so the delivered engine, XeLaTeX with its default Latin Modern text font, reports no missing character.
	\begin{theorem}
		\label{fracdyntm}
		The solution $\upsilon \in C([0,\infty),(0,\infty)) \cap D_0$ to the problem
		\begin{align*}
			\begin{cases}
				\displaystyle 	\dot{\upsilon}(t,x)= \upsilon^{\prime \prime}(t,x) & t>0,\; x \in (0,\infty),\\
				\displaystyle 	\eta\, \mathfrak{D}_t^\Psi \upsilon(t,x) = -\mathbf{D}_{x-}^\Phi \upsilon(t,x), \; \eta \geq 0    & t>0,\; x=0,\\
				\displaystyle 	\upsilon(0,x)=f(x), \; f \in C[0,\infty) \cap L^\infty(0, \infty)  &x\in [0, \infty)
			\end{cases}
		\end{align*}
		has the probabilistic representation
		\begin{align}
			\label{NLstickyBpallino}
			\upsilon(t,x)=\mathbf{E}_x\left[f(B^\bullet \circ S_t^{-1})\right].
		\end{align}
	\end{theorem}

		\begin{align}
			\label{upsilonSol}
			\upsilon(t,x)=Q_t^D f(x)+\int_0^t \frac{x}{\tau} g(\tau,x) \upsilon(t-\tau,0) d\tau
		\end{align}

		\begin{align}
			\label{timeBondCond}
			\int_0^\infty e^{-\lambda t} \eta\mathfrak{D}_t^\Psi \upsilon(t,x) dt \bigg\vert_{x=0}= \eta \Psi(\lambda) R_\lambda f(0) - \eta \frac{\Psi(\lambda)}{\lambda} f(0), \quad \lambda>0.
		\end{align}

		\begin{align}
			I_1^\Psi+I_2^\Psi
			%	&=\frac{\int_0^\infty R_\lambda^D f(y) \, \Pi^\Phi(dy)}{\eta \Psi(\lambda) +\Phi(\sqrt \lambda)} + \frac{\eta}{\lambda} \Psi(\lambda) f(0) \frac{1}{\eta \Psi(\lambda)+\Phi(\sqrt \lambda)}\\
			&=\frac{1}{\eta \Psi(\lambda) +\Phi(\sqrt \lambda)} \left( \int_0^\infty R_\lambda^D f(y) \, \Pi^\Phi(dy) + \frac{\eta}{\lambda} \Psi(\lambda)f(0) \right)
			\label{I1piuI2}
		\end{align}

		\begin{align}
			\label{resolvent2}
			\mathbf{E}_\mathsf{x}\left[e^{-\lambda \mathcal{T}}\right] \mathbf{E}_v\left[\int_{0}^{\infty} e^{-\lambda t} f(\mathcal{X}^\bullet_t) dt\right] =  e^{-x \sqrt{\lambda}}  \sum_{e \in \mathcal{E}} \rho_e \mathbf{E}_{(e,0)}  \left[\int_0^\infty e^{-\lambda t} f(e, B^\bullet_t) dt\right].
		\end{align}

		\begin{align}
			\label{marchaud-graphs}
			0&=- \sum_{e \in \mathcal{E}} \rho_e \mathbf{D}_{x-}^\Phi \mathcal{R}_\lambda f(e,x)\big\vert_{x=0}\\
			& =\sum_{e \in \mathcal{E}} \rho_e \int_0^\infty \mathcal{R}_\lambda^D f(e,y)\, \Pi^\Phi(dy) - \Phi(\sqrt{\lambda})\, \sum_{e \in \mathcal{E}} \rho_e {\mathcal{R}}_\lambda f(e,0),
		\end{align}

		\begin{theorem}
			\label{thm:Xpallinodelay}
			The probabilistic representation of the solution $u \in D_H \cap D_L$ of the problem
			\begin{align*}
				\begin{cases}
					\dot{u}(t,\mathsf{x})= \Delta u(t,\mathsf{x}) \quad &t>0, \mathsf{x} \in \mathcal{G}\setminus \{v\}, \\
					\eta \mathfrak{D}_t^\Psi u(t,v)+ \sum_{e \in \mathcal{E}} \rho_e\mathbf{D}^\Phi_{x-} u_e(t,x)\big\vert_{x=0}=0 \quad    &t>0,\\
					u(0,\mathsf{x})=f(\mathsf{x}) \quad &\mathsf{x} \in \mathcal{G},
				\end{cases}
			\end{align*}
			with $f \in C(\mathcal{G}) \cap L^\infty(\mathcal{G})$, $\eta \geq 0$  and $\mathsf{x}=(l,x)$ is
			\begin{align*}
				u(t,\mathsf{x})=\mathbf{E}_\mathsf{x}[f(\mathcal{X}^\bullet \circ \mathcal{S}^{-1}_t)].
			\end{align*}
		\end{theorem}

		\begin{proof}
			The proof of this theorem follows that of Theorem \ref{fracdyntm}, where the key observation is that once the edge is specified, the dynamics on $(0,\infty)$ remain the same. Let us examine the principal steps.
			
			We write the $\lambda-$potential, for $\lambda >0$,
			\begin{align*}
				\mathcal{R}_\lambda f(\mathsf{x})=\int_{0}^{\infty} e^{-\lambda t} u(t,\mathsf{x}) dt =\int_{0}^{\infty} e^{-\lambda t} \mathbf{E}_\mathsf{x}[f(\mathcal{X}^\bullet \circ \mathcal{S}^{-1}_t)] dt.
			\end{align*}
			Since $\mathsf{x}=(l,x)$, before hitting the vertex the motion is a killed Brownian motion on the edge $e$. Then, with the same arguments that led us to write \eqref{upsilonSol},  and once a vertex is reached, we choose the edge as in \eqref{resolvent2}, we have, for $\lambda>0$,
			\begin{align}
				\label{resolvent-graph}
				\mathcal{R}_\lambda f(\mathsf{x})=\mathcal{R}_\lambda^D f(\mathsf{x})+ e^{-x \sqrt{\lambda}}  \sum_{e \in \mathcal{E}} \rho_e \mathbf{E}_{(e,0)}  \left[\int_0^\infty e^{-\lambda t} f(e, B^\bullet\circ {S}^{-1}_t) dt\right],
			\end{align}
			where ${S}^{-1}$ is the same random time introduced in Theorem \ref{fracdyntm}, this is possible since the local time $\ell$ at the vertex coincides with $\gamma$.
			
			
			As previously noted, the heat equation holds, and we only need to verify the conditions at the vertex. 
			As mentioned earlier, once the edge is specified, the $\lambda-$potential to be computed remains that of the process $B^\bullet\circ {S}^{-1}$ on $(0,\infty)$, which is already determined in \eqref{I1piuI2} through the sum of $I_1^\Psi+ I_2^\Psi$. Then, we have
			\begin{align*}
				&\mathbf{E}_{(e,0)}  \left[\int_0^\infty e^{-\lambda t} f(e, B^\bullet\circ {S}^{-1}_t) dt\right]\\&=\frac{1}{\eta \Psi(\lambda) +\Phi(\sqrt \lambda)} \left( \int_0^\infty R_\lambda^D f(e,y) \, \Pi^\Phi(dy) + \frac{\eta}{\lambda} \Psi(\lambda)f(e,0) \right)
			\end{align*}
			and the $\lambda-$potential at the vertex, by using \eqref{resolvent-graph}, is
			\begin{align*}
				\mathcal{R}_\lambda f(v)= \sum_{e\in \mathcal{E}} \rho_e \left(\frac{1}{\eta \Psi(\lambda) +\Phi(\sqrt \lambda)} \left( \int_0^\infty R_\lambda^D f(e,y) \, \Pi^\Phi(dy) + \frac{\eta}{\lambda} \Psi(\lambda)f(e,0) \right)\right).
			\end{align*}
			We now verify the Laplace transforms for the conditions at the vertex. For the Caputo-D\v{z}rba\v{s}jan (type) derivative, as in \eqref{timeBondCond}, we get
			\begin{align}
				\label{caputo-graphs-bc}
				\int_0^\infty e^{-\lambda t} \eta\mathfrak{D}_t^\Psi u(t,v) dt = \eta \Psi(\lambda) \mathcal{R}_\lambda f(v) - \eta \frac{\Psi(\lambda)}{\lambda} f(v), \quad \lambda>0,
			\end{align}
			where we emphasize that, in exploring solutions within space $D_H$, the function at the vertex is continuous. For the Marchaud (type) derivative, from \eqref{marchaud-graphs}, we have
			\begin{align}
				\label{marchaud-graphs-BC}
				\sum_{e \in \mathcal{E}} \rho_e \mathbf{D}_{x-}^\Phi \mathcal{R}_\lambda f(e,x)\big\vert_{x=0} =-\sum_{e \in \mathcal{E}} \rho_e \int_0^\infty \mathcal{R}_\lambda^D f(e,y)\, \Pi^\Phi(dy) + \Phi(\sqrt{\lambda})\, \mathcal{R}_\lambda f(v).
			\end{align}
			By combining \eqref{caputo-graphs-bc} and \eqref{marchaud-graphs-BC}, the vertex conditions applied to the $\lambda-$potential become
			\begin{align*}
				\eta \Psi(\lambda) \mathcal{R}_\lambda f(v) - \eta \frac{\Psi(\lambda)}{\lambda} f(v) - \sum_{e \in \mathcal{E}} \rho_e \int_0^\infty \mathcal{R}_\lambda^D f(e,y)\, \Pi^\Phi(dy) + \Phi(\sqrt{\lambda})\, \mathcal{R}_\lambda f(v)=0
			\end{align*}
			which are satisfied by
			\begin{align*}
				\mathcal{R}_\lambda f(v)= \sum_{e\in \mathcal{E}} \rho_e \left(\frac{1}{\eta \Psi(\lambda) +\Phi(\sqrt \lambda)} \left( \int_0^\infty R_\lambda^D f(e,y) \, \Pi^\Phi(dy) + \frac{\eta}{\lambda} \Psi(\lambda)f(e,0) \right)\right),
			\end{align*}
			that is the $\lambda-$potential at the vertex of the process $\mathcal{X}^\bullet \circ \mathcal{S}^{-1}$. This concludes that $u(t,\mathsf{x})=\mathbf{E}_\mathsf{x}[f(\mathcal{X}^\bullet \circ \mathcal{S}^{-1}_t)]$ is a solution to the problem under consideration.
		\end{proof}

\end{document}
